% Partie : set-theory
% Chapitre : ordinals
% Section : basic

\documentclass[../../../include/open-logic-section]{subfiles}

\begin{document}

\olfileid[fr]{sth}{ordinals}{basic}
\olsection{Propriétés élémentaires des ordinaux}

Nous avons constaté que les premiers ordinaux sont les nombres
naturels. La principale raison de développer une théorie des ordinaux
est d'étendre le principe de récurrence valable pour les nombres
naturels. Nous y parviendrons par une série de résultats élémentaires.

\begin{lem}\ollabel{ordmemberord}
Tout !!{element} d'un ordinal est un ordinal.
\end{lem}

\begin{proof}
Soit $\alpha$ un ordinal, avec $b \in \alpha$. Puisque $\alpha$
est transitif, $b \subseteq \alpha$. La relation $\in$ est donc un
bon ordre sur $b$, puisqu'elle est un bon ordre sur $\alpha$.

Pour voir que $b$ est transitif, supposons $x \in c \in b$. Comme
$b \subseteq \alpha$, nous avons $c \in \alpha$. La transitivité de
$\alpha$ donne alors $c \subseteq \alpha$, donc $x \in \alpha$.
Ainsi, $x, c, b \in \alpha$. Or $\in$ est un bon ordre sur $\alpha$
et, par la \olref[sth][ordinals][wo]{wo:strictorder}, une relation
transitive sur $\alpha$. De $x \in c \in b$, il suit donc que
$x \in b$. En général, $c \subseteq b$.
\end{proof}

\begin{cor}\ollabel{ordissetofsmallerord}
$\alpha = \Setabs{\beta \in \alpha}{\beta \text{ est un ordinal}}$
pour tout ordinal~$\alpha$.
\end{cor}

\begin{proof}
Cela découle immédiatement du \olref{ordmemberord}.
\end{proof}

L'idée générale des deux résultats principaux qui suivent,
\olref{ordinductionschema} et \olref{ordtrichotomy}, est que les
ordinaux eux-mêmes sont bien ordonnés par l'appartenance :

\begin{thm}[Récurrence transfinie]\ollabel{ordinductionschema}
Pour toute formule $\phi(x)$ :
\[
	\text{si }\exists \alpha \phi(\alpha)\text{, alors }\exists \alpha(\phi(\alpha)
	\land  (\forall \beta \in \alpha) \lnot \phi(\beta))
\]
où les quantificateurs affichés portent implicitement sur les seuls
ordinaux.
\end{thm}

\begin{proof}
%\olref{ordinductionschema1}.
Supposons que $\phi(\alpha)$ soit vraie pour un ordinal $\alpha$. Si
$(\forall \beta \in \alpha) \lnot \phi(\beta)$, c'est terminé.
Sinon, parmi les éléments de $\alpha$ qui satisfont $\phi$,
il en existe un plus petit pour $\in$, puisque $\alpha$ est un ordinal.
Cet élément est lui-même un ordinal par le
\olref{ordmemberord}.
% Réaudit : la source choisit déjà le plus petit témoin de phi dans alpha.
% L'expansion explicite par séparation, correcte mais facultative, est
% conservée dans le dossier de propositions et retirée du corps traduit.
%
%
%
%\olref{ordinductionschema2}. Supposons qu'il existe un ordinal
%\gamma tel que \lnot\phi(\gamma). Alors, par
%\olref{ordinductionschema1}, il existe un ordinal \in-minimal
%\alpha tel que \lnot\phi(\alpha). Donc (\forall \beta <
%\alpha) \phi(\beta), ce qui rend fausse la prémisse de l'implication.
\end{proof}
\noindent
Notons que nous pouvons aussi exprimer \olref{ordinductionschema}
par le schéma :
\[
\text{si }\forall \alpha((\forall \beta \in \alpha)\phi(\beta) \lif
\phi(\alpha))\text{, alors }\forall \alpha\phi(\alpha)
\]
Il suffit d'appliquer le \olref{ordinductionschema} à
$\lnot\phi(\alpha)$, puis d'effectuer quelques transformations
logiques élémentaires.

\begin{thm}[Trichotomie]\ollabel{ordtrichotomy}
$\alpha \in \beta \lor \alpha = \beta \lor \beta \in \alpha$, pour
tous ordinaux $\alpha$ et $\beta$.
\end{thm}

\begin{proof}
La démonstration procède par double récurrence, c'est-à-dire en
appliquant deux fois \olref{ordinductionschema}. Disons que $x$ est
\emph{comparable} à $y$ si et seulement si
$x \in y \lor x = y \lor y \in x$.

Comme hypothèse de la première récurrence, supposons que tout ordinal
appartenant à~$\alpha$ soit comparable à \emph{tout} ordinal. Comme
hypothèse de la seconde, supposons que $\alpha$ soit comparable à
tout ordinal appartenant à~$\beta$. Nous allons montrer que $\alpha$
est comparable à~$\beta$. Par récurrence sur~$\beta$, il s'ensuivra
que $\alpha$ est comparable à tout ordinal ; puis, par récurrence
sur~$\alpha$, que \emph{tout} ordinal est comparable à \emph{tout}
ordinal, comme voulu. Il suffit de supposer $\alpha \notin \beta$ et
$\beta \notin \alpha$, puis de montrer que $\alpha = \beta$.

Pour établir $\alpha \subseteq \beta$, fixons $\gamma \in \alpha$ ;
c'est un ordinal par le \olref{ordmemberord}. La première hypothèse de
récurrence assure donc que $\gamma$ est comparable à $\beta$. Mais,
si $\gamma = \beta$ ou $\beta \in \gamma$, alors
$\beta \in \alpha$ (en invoquant, au besoin, la transitivité de
$\alpha$), contrairement à notre hypothèse. Par conséquent,
$\gamma \in \beta$. Ainsi, $\alpha \subseteq \beta$.

Un raisonnement exactement analogue, fondé sur la seconde hypothèse
de récurrence, établit $\beta \subseteq \alpha$. Donc
$\alpha = \beta$.
\end{proof}\noindent Nous écrirons donc parfois $\alpha <\beta$ au
lieu de $\alpha \in \beta$, puisque $\in$ se comporte comme une
relation d'ordre. Ce choix ne repose sur aucune raison profonde :
cette notation est familière, et il est plus commode d'écrire
$\alpha \leq \beta$ que $\alpha \in \beta \lor \alpha = \beta$.
\footnote{Nous pourrions écrire $\alpha
\mathrel{\underline{\in}} \beta$, mais cette notation serait tout à
fait inhabituelle.}

Voici deux conséquences rapides des derniers résultats, dont la
première met en œuvre notre nouvelle notation :

\begin{cor}\ollabel{ordordered}
Si $\exists \alpha\phi(\alpha)$, alors $\exists \alpha(\phi(\alpha)
\land \forall \beta(\phi(\beta) \lif \alpha \leq \beta))$.
De plus, pour tous ordinaux $\alpha, \beta, \gamma$, on a
$\alpha \notin \alpha$ et
$\alpha \in \beta \in \gamma \lif \alpha \in \gamma$.
\end{cor}

\begin{proof}
Comme dans la \olref[wo]{wo:strictorder}.
\end{proof}

\begin{prob}
Complétez le «~raisonnement exactement analogue~» dans la
démonstration du \olref[sth][ordinals][basic]{ordtrichotomy}.
\end{prob}

\begin{cor}\ollabel{corordtransitiveord}
$A$ est un ordinal si et seulement si $A$ est un ensemble transitif
d'ordinaux.
\end{cor}

\begin{proof}
\emph{De gauche à droite.} Par le \olref{ordmemberord}.
\emph{De droite à gauche.} Si $A$ est un ensemble transitif
d'ordinaux, alors $\in$ est un bon ordre sur $A$ par le
\olref{ordinductionschema} et le \olref{ordtrichotomy}.
\end{proof}

Nous avons présenté le \olref{ordinductionschema} et le
\olref{ordtrichotomy} comme affirmant que $\in$ est un bon ordre sur
les ordinaux. Il faut pourtant être \emph{très prudent} avec ce genre
d'affirmation, à cause du résultat suivant :

\begin{thm}[Paradoxe de Burali-Forti]\ollabel{buraliforti}
Il n'existe pas d'ensemble de tous les ordinaux.
\end{thm}

\begin{proof}
Par l'absurde, supposons que $O$ soit l'ensemble de tous les
ordinaux. Si $\alpha \in \beta \in O$, alors $\alpha$ est un ordinal
par le \olref{ordmemberord}, donc $\alpha \in O$. Ainsi, $O$ est
transitif, et c'est donc un ordinal par \olref{corordtransitiveord}.
Il s'ensuit que $O \in O$, ce qui contredit \olref{ordordered}.
\end{proof}
\noindent
Ce résultat porte le nom de \citeauthor{Burali-Forti1897}. Mais c'est
Cantor qui, en 1899, dans une lettre à Dedekind, a le premier vu
clairement la \emph{contradiction} qu'entraîne l'hypothèse d'un
ensemble de tous les ordinaux. Comme l'explique van Heijenoort :
\begin{quote}
  Burali-Forti lui-même considérait que la contradiction établissait,
  par \emph{réduction à l'absurde}, que l'ordre naturel des ordinaux
  n'est qu'un ordre partiel.
  \citep[p.~105]{Heijenoort1967}
\end{quote}
Si l'on met de côté l'erreur de Burali-Forti, on peut résumer ainsi
ce qui précède : les ordinaux sont des ensembles dont chacun est bien
ordonné par l'appartenance, et qui sont collectivement bien ordonnés
par elle, sans constituer collectivement un ensemble.

Pour terminer, voici encore quelques propriétés élémentaires des
ordinaux qui découlent du \olref{ordinductionschema} et du
\olref{ordtrichotomy}.

\begin{prop}
Toute suite strictement décroissante d'ordinaux est finie.
\end{prop}

\begin{proof}
Une suite infinie strictement décroissante d'ordinaux
$\alpha_0 > \alpha_1 > \alpha_2 > \ldots$ ne possède aucun
élément minimal pour $<$, ce qui contredit
\olref{ordinductionschema}.
\end{proof}

\begin{prop}\ollabel{ordinalsaresubsets}
$\alpha \subseteq \beta \lor \beta \subseteq \alpha$, pour tous
ordinaux $\alpha, \beta$.
\end{prop}

\begin{proof}
Si $\alpha \in \beta$, alors $\alpha \subseteq \beta$ puisque
$\beta$ est transitif. De même, si $\beta \in \alpha$, alors
$\beta \subseteq \alpha$. Enfin, si $\alpha = \beta$, alors
$\alpha \subseteq \beta$ et $\beta \subseteq \alpha$.
\olref{ordtrichotomy} permet donc de conclure.
\end{proof}

\begin{prop}\ollabel{ordisoidentity}
$\alpha = \beta$ si et seulement si $\ordeq{\alpha}{\beta}$, pour
tous ordinaux $\alpha, \beta$.
\end{prop}

\begin{proof}
Les ordinaux sont des bons ordres ; le résultat découle donc
immédiatement de la trichotomie (\olref[basic]{ordtrichotomy}) et du
\olref[iso]{wellordnotinitial}.
\end{proof}

\begin{prob}\ollabel{probunionordinalsordinal}
Démontrez que, si tout élément de $X$ est un ordinal, alors
$\bigcup X$ est un ordinal.
\end{prob}

\end{document}
